\scanpage{128}
\mspage{115}
se ramène aux deux relations
\[
\begin{aligned}
&\underbrace{\begin{array}[t]{c@{\,}c@{\,}c}
 \underline{(\beta_0\widetilde i)r}&=&\underline{\widetilde\beta_0r},\\
 \Vert&&\Vert\\
 \beta_0(\widetilde i r)&&p_0\alpha_0\\
 \Vert&&\\
 \beta_0i_0&&
\end{array}}_{\substack{\text{OK par comm.}\\\text{du carré du bas}\\\text{de (6.3.5)}}}
\quad
\underbrace{\begin{array}[t]{c@{\,}c@{\,}c}
 \underline{(\beta_0\widetilde i)s}&=&\underline{\widetilde\beta_0s}\\
 \Vert&&\Vert\\
 \beta_0(\widetilde i s)&&\lambda_Y\beta_{I'}\\
 \Vert&&\\
 \beta_0\lambda_B&&
\end{array}}_{\substack{\text{OK. par comm.}\\\text{du carré antérieur}\\\text{de (6.3.6).}}}
\end{aligned}
\]
Donc il reste à vérifier \deleted{les} la \inserted{deuxième} \unsure{relations|relation} (6.3.12)
\[
 \deleted{\widetilde h_0 r=\alpha_0,}\quad
 p_0\widetilde h_0=\widetilde\beta_0,\quad
 \deleted{\unsure{\lambda_X h'_{I'}}\ \widetilde h_0 s=\lambda_X h'_{I'}}
\]
ce qui se ramène à
\[
\begin{aligned}
&\begin{array}[t]{c@{\,}c@{\,}c}
 \underline{(p_0\widetilde h_0)r}&=&\widetilde\beta_0 r,\\
 \Vert&&\\
 p_0(\underset{\alpha_0}{\underline{\widetilde h_0r}})&&\\
 \underline{\phantom{p_0\widetilde h_0r}}&&\\[-2pt]
 p_0\alpha_0&&\\
 \multicolumn{3}{c}{\underbrace{\hphantom{(p_0\widetilde h_0)r=\widetilde\beta_0r}}_{\text{OK par (6.3.11)}}}
\end{array}
\qquad
\begin{array}[t]{c@{\,}c@{\,}c}
 \underline{(p_0\widetilde h_0)s}&=&\underline{\widetilde\beta_0s}\\
 p_0(\underset{\lambda_Xh'_{I'}}{\underline{\widetilde h_0s}})&&\lambda_Y\beta_{I'}\\
 \underline{\phantom{(p_0\lambda_X)h'_{I'}}}&&\\[-2pt]
 (\underset{\lambda_Yp_{I'}}{\underline{p_0\lambda_X}})h'_{I'}&&\\
 \underline{\phantom{\lambda_Y(p_{I'}h'_{I'})}}&&\\[-2pt]
 \deleted{\unsure{p_I|\lambda_I}}\unsure{\lambda_Y}(\underset{\unsure{p_{I'}|\beta_{I'}}}{\underline{p_{I'}h'_{I'}}})&&\\
 \multicolumn{3}{c}{\underbrace{\hphantom{(p_0\widetilde h_0)s=\widetilde\beta_0s}}_{\text{OK}}}
\end{array}
\end{aligned}
\]

\deleted{Je dis}
Ceci posé, je dis que les contraintes
(6.3.7), (6.3.8) sur $h_0$ reviennent à la
commutativité des deux triangles dans
le carré

\diagram{\((6.3.14)\)\qquad
\begin{tikzcd}[ampersand replacement=\&]
 \widetilde B\arrow[r,"{\widetilde h_0}"]\arrow[d,"{\widetilde i\in}"',"\Phi"] \& X_0\arrow[d,"{p_0\in\Psi}"]\\
 B_0\arrow[r,"{\beta_0}"]\arrow[ur,dotted,"{h_0}"] \& Y_0,
\end{tikzcd}}
% archive sheet number: 127

